% 可直接粘贴到 Overleaf 主文档中的第三部分
% 推荐编译器：XeLaTeX；图片放在 figures/ 目录

\section{RadioDiff 基础实验复现}
\label{sec:radiodiff-reproduction}

本节完成 RadioDiff 的基础复现。实验使用 RadioMapSeer 数据集、本地预训练权重
\texttt{model-irt.pt} 和单张 NVIDIA RTX 5060 Laptop GPU。每组实验只推理测试集中的前
5 个样本，样本名称为 \texttt{289\_0} 至 \texttt{289\_4}，不重新训练 RadioDiff。

\subsection{问题定义与核心假设}

给定城市建筑物、车辆和发射机的空间条件，目标是直接生成全空间 pathloss radio map
（RM），避免逐点现场测量或高成本射线追踪。本文验证三个假设：
（1）建筑物、车辆和发射机条件包含决定 RM 大尺度结构的主要信息；
（2）RM 中的遮挡、多峰和局部传播结构具有条件生成分布特征，扩散模型适合对其建模；
（3）在 checkpoint 对应的数据分布内，模型能够从条件输入恢复测试场景的传播趋势。

\subsection{实现方法}

RadioDiff 采用两阶段 latent diffusion。第一阶段的 AutoencoderKL 将单通道 RM 压缩为
三通道 latent，第二阶段的条件 U-Net 在 latent 空间中执行去噪。条件分支使用 Swin
风格的窗口特征和 Fourier 特征，分别用于提取建筑物边界、车辆变化、发射机位置及
更大尺度的空间变化。推理时使用滑动窗口，对重叠窗口的结果进行拼接。

本地调用链为：
\begin{verbatim}
radio_sample.yaml -> RadioUNet_c -> AutoencoderKL
 -> conditional U-Net -> LatentDiffusion
 -> slide_sample -> PNG and metrics.json
\end{verbatim}

配置文件中的 \texttt{data.name=radio} 对应 \texttt{RadioUNet\_c}，其默认目标为
\texttt{gain/carsDPM}，条件输入为建筑物、发射机和车辆三通道。checkpoint 中包含
1536 个模型参数项，其中 372 个属于 \texttt{first\_stage\_model}，因此本次推理
同时恢复了 VAE 参数。

\subsection{实验设置}

模型图像尺寸为 $320\times320$，batch size 为 1，采样类型为滑窗采样，stride 为
$(256,256)$。每次实验推理 5 个样本，并保存生成 RM、解析后的配置、运行信息和指标。
主要评价指标为 NMSE、RMSE、SSIM 和 PSNR。其中 NMSE、RMSE 越低越好，SSIM、PSNR 越高越好。

\begin{table}[htbp]
  \centering
  \caption{RadioDiff 推理对照实验设置}
  \label{tab:exp-settings}
  \begin{tabular}{llll}
    \toprule
    实验 & 随机种子 & 扩散步数 & 目的 \\
    \midrule
    baseline\_steps\_5 & 42 & 5 & 基线 \\
    baseline\_steps\_20 & 42 & 20 & 扩散步数对照 \\
    baseline\_steps\_50 & 42 & 50 & 扩散步数对照 \\
    seed\_123 & 123 & 5 & 随机性对照 \\
    \bottomrule
  \end{tabular}
\end{table}

\subsection{关键实验的预先预测}

在运行实验前，先记录以下预测，运行完成后再与真实结果比较。

\paragraph{实验 E1：扩散步数。}
固定随机种子为 42，比较 5、20 和 50 步。预期 5 步已经可以生成合理 RM；增加步数
可能在一定范围内改善结果，但继续增加不一定改善，并且推理时间会明显增加。原因是
当前 checkpoint、采样器和训练过程对应特定的离散化设置，采样步数不是越大越接近最优。

\paragraph{实验 E2：随机种子。}
固定扩散步数为 5，比较 seed=42 和 seed=123。预期传播结构基本保持，局部细节和误差
会有小幅变化。建筑物和发射机条件决定主要结构，初始噪声更多影响局部细节。如果两组
结果完全相同，则需要检查采样器是否在内部覆盖了配置中的随机种子。

\subsection{定量结果}

\begin{table}[htbp]
  \centering
  \caption{五个测试样本的平均指标}
  \label{tab:quant-results}
  \begin{tabular}{lrrrrrr}
    \toprule
    实验 & seed & steps & NMSE$\downarrow$ & RMSE$\downarrow$ & SSIM$\uparrow$ & PSNR$\uparrow$ \\
    \midrule
    baseline\_steps\_5 & 42 & 5 & 0.010148 & 0.039451 & 0.942289 & 28.1332 \\
    baseline\_steps\_20 & 42 & 20 & 0.010325 & 0.039745 & 0.937611 & 28.0739 \\
    baseline\_steps\_50 & 42 & 50 & 0.010602 & 0.040294 & 0.936071 & 27.9523 \\
    seed\_123 & 123 & 5 & 0.010343 & 0.039785 & 0.941220 & 28.0642 \\
    \bottomrule
  \end{tabular}
\end{table}

5 步 seed=42 的综合指标最好。20 步和 50 步没有继续改善，反而使 NMSE、RMSE 增大，
SSIM 和 PSNR 降低。seed=123 相比 seed=42 只产生小幅变化，说明本组样本的主要传播
结构由条件输入控制，随机种子主要影响局部生成细节。

\subsection{可视化结果}

图~\ref{fig:input-condition} 展示样本 289\_3 的建筑物、发射机和车辆三通道条件。
图~\ref{fig:qualitative-best} 是逐样本 NMSE 最低的样本 289\_3，图~\ref{fig:qualitative-worst}
是逐样本 NMSE 最高的样本 289\_2。图~\ref{fig:steps-comparison} 对比同一样本在不同
扩散步数和随机种子下的输出。

\begin{figure}[htbp]
  \centering
  \includegraphics[width=0.98\linewidth]{figures/input_condition.png}
  \caption{样本 289\_3 的条件输入：建筑物、发射机和车辆。}
  \label{fig:input-condition}
\end{figure}

\begin{figure}[htbp]
  \centering
  \includegraphics[width=0.88\linewidth]{figures/qualitative_best.png}
  \caption{较好结果：样本 289\_3 的真实 RM、预测 RM 和绝对误差。}
  \label{fig:qualitative-best}
\end{figure}

\begin{figure}[htbp]
  \centering
  \includegraphics[width=0.88\linewidth]{figures/qualitative_worst.png}
  \caption{较差结果：样本 289\_2 的真实 RM、预测 RM 和绝对误差。}
  \label{fig:qualitative-worst}
\end{figure}

\begin{figure}[htbp]
  \centering
  \includegraphics[width=\linewidth]{figures/steps_comparison.png}
  \caption{样本 289\_2 在不同扩散步数和随机种子下的输出对比。}
  \label{fig:steps-comparison}
\end{figure}

\subsection{好样本与差样本分析}

本文根据逐样本 NMSE 对五个测试样本进行排序，选择样本 289\_3 作为较好结果，选择样本
289\_2 作为较差结果。解释部分预留为空白，待根据可视化结果和个人理解补充。

\begin{table}[htbp]
  \centering
  \small
  \renewcommand{\arraystretch}{1.35}
  \setlength{\tabcolsep}{5pt}
  \caption{好样本与差样本对比}
  \label{tab:good-bad-samples}
  \begin{tabularx}{\linewidth}{
    >{\centering\arraybackslash}p{1.1cm}
    >{\centering\arraybackslash}p{1.1cm}
    >{\raggedright\arraybackslash}p{3.0cm}
    L
  }
    \toprule
    样本 & 表现 & 定量指标 & 解释 \\
    \midrule
    289\_3 & 最好
      & NMSE=0.008549；RMSE=0.036897；SSIM=0.921124
      & \rule{0pt}{2.5cm} \\
    289\_2 & 最差
      & NMSE=0.018559；RMSE=0.050659；SSIM=0.887862
      & \rule{0pt}{2.5cm} \\
    \bottomrule
  \end{tabularx}
\end{table}

\noindent\textbf{补充分析：}

\vspace{3cm}

\subsection{挑战、修正与结果解释}

本次最大挑战是确认实验变量确实生效。第一次修改 seed 后，seed=42 和 seed=123 的
PNG 文件哈希完全一致。沿采样调用链检查后发现，\texttt{ddm\_const\_sde.py} 的采样
函数内部固定执行了 \texttt{torch.manual\_seed(0)}、\texttt{numpy.seed(0)} 和
\texttt{random.seed(0)}，覆盖了配置值。

解决方法是增加 \texttt{sampler.seed} 配置，将其传入 \texttt{model\_cfg.seed}，并在
采样函数内读取该值。修正后重新运行全部关键对照，seed=123 的 NMSE 为 0.010343，
与 seed=42 的 0.010148 存在小幅差异，证明随机种子变量已经真正生效。

如果继续三天，第一优先级是扩大并分层评估更多城市、发射机位置和车辆状态，分别统计
建筑边界、发射机附近、遮挡区和远距离区域的误差。第二优先级是严格验证 loader、
gain 类型和 checkpoint 的配对关系，避免将 DPM、IRT4、DPMCAR 等不同条件混用。第三
优先级是加入发射机附近衰减、建筑物穿透惩罚和遮挡区域误差等物理一致性检查；相比盲目
增加扩散步数，这些改动更可能提升实际泛化能力。

\subsection{本节结论}

本地实验满足基础复现的最低要求：预训练模型成功运行，得到四组独立 RM 结果，报告了
多项定量指标，分析了较好和较差样本，并改变了扩散步数和随机种子两个推理条件。结论
限定在当前 checkpoint、数据 loader 和五个测试样本上；若要得到更稳定的性能判断，还
需要扩大测试集并进行跨场景评估。
